\documentclass[tikz,border=8pt]{standalone}
\usepackage{tikz}
\usetikzlibrary{calc}
\usepackage{amsmath}
\usepackage{pgfplots}
\pgfplotsset{compat=1.18}
\usepackage{ctex}

\begin{document}
\begin{tikzpicture}

%% ── 左图：X̄ 抽样分布，n=5 vs n=30 ─────────────────────────────
\begin{axis}[
  name=ax1,
  width=7.0cm, height=5.5cm,
  xmin=-3.5, xmax=3.5,
  ymin=0, ymax=1.32,
  xlabel={$\bar{x}$},
  ylabel={$f(\bar{x})$},
  axis lines=left,
  xtick={-3,-2,-1,0,1,2,3},
  xticklabel style={font=\scriptsize},
  yticklabel style={font=\scriptsize},
  ytick={0,0.2,0.4,0.6,0.8,1.0,1.2},
  grid=major, grid style={gray!20},
  clip=false,
  legend style={font=\scriptsize, at={(0.98,0.95)}, anchor=north east}
]
  % n=5: N(0, 1/5)，sd=sqrt(1/5)≈0.447
  \addplot[blue!70!black, thick, domain=-3.5:3.5, samples=150]
    {exp(-x^2/(2*0.2))/sqrt(2*3.14159*0.2)};
  \addlegendentry{$n=5$，$N(\mu,\sigma^2/5)$}

  % n=30: N(0, 1/30)，sd=sqrt(1/30)≈0.183
  \addplot[red!70!black, thick, dashed, domain=-3.5:3.5, samples=150]
    {exp(-x^2/(2*0.0333))/sqrt(2*3.14159*0.0333)};
  \addlegendentry{$n=30$，$N(\mu,\sigma^2/30)$}

  % μ 参考线
  \addplot[gray, dashed] coordinates {(0,0) (0,1.25)};
  \node[font=\scriptsize, below] at (axis cs:0,0) {$\mu$};

  \node[font=\scriptsize, align=left, fill=white, fill opacity=0.85, text opacity=1,
        rounded corners=2pt] at (axis cs:-2.2,1.15)
    {$n$ 增大 $\Rightarrow$\\方差 $\sigma^2/n$ 减小};
\end{axis}

%% ── 右图：(n-1)S²/σ² 的 χ² 分布，df=5/10/30 ─────────────────
% χ² 近似值（手工计算 PDF 关键点）
% df=5:  peak≈0.175 at x≈3;   df=10: peak≈0.096 at x≈8;  df=30: peak≈0.043 at x≈28
\begin{axis}[
  name=ax2,
  at={(ax1.east)}, anchor=west, xshift=0.6cm,
  width=7.0cm, height=5.5cm,
  xmin=0, xmax=55,
  ymin=0, ymax=0.22,
  xlabel={$\chi^2$},
  ylabel={$f(\chi^2)$},
  axis lines=left,
  xtick={0,10,20,30,40,50},
  xticklabel style={font=\scriptsize},
  yticklabel style={font=\scriptsize},
  ytick={0,0.05,0.10,0.15,0.20},
  grid=major, grid style={gray!20},
  clip=false,
  legend style={font=\scriptsize, at={(0.98,0.95)}, anchor=north east}
]
  % χ²(5) 写死关键点（PDF 用 spline 拟合）
  \addplot[blue!70!black, thick, smooth] coordinates {
    (0.2,0.044)(0.5,0.095)(1.0,0.152)(1.5,0.170)(2.0,0.174)(2.5,0.175)
    (3.0,0.172)(3.5,0.165)(4.0,0.157)(4.5,0.148)(5.0,0.138)(6.0,0.119)
    (7.0,0.100)(8.0,0.082)(9.0,0.066)(10.0,0.053)(12.0,0.032)(14.0,0.018)
    (16.0,0.010)(18.0,0.005)(20.0,0.002)(22.0,0.001)
  };
  \addlegendentry{$df=5$}

  % χ²(10)
  \addplot[red!70!black, thick, dashed, smooth] coordinates {
    (0.5,0.006)(1.0,0.017)(2.0,0.038)(3.0,0.057)(4.0,0.071)(5.0,0.081)
    (6.0,0.087)(7.0,0.090)(8.0,0.092)(9.0,0.091)(10.0,0.088)(11.0,0.084)
    (12.0,0.079)(13.0,0.072)(14.0,0.065)(15.0,0.058)(16.0,0.051)(18.0,0.038)
    (20.0,0.027)(22.0,0.019)(25.0,0.010)(28.0,0.005)(32.0,0.002)
  };
  \addlegendentry{$df=10$}

  % χ²(30)
  \addplot[green!60!black, thick, dotted, smooth] coordinates {
    (10,0.003)(12,0.008)(14,0.015)(16,0.022)(18,0.029)(20,0.035)(22,0.040)
    (24,0.043)(26,0.045)(28,0.045)(30,0.043)(32,0.040)(34,0.036)(36,0.032)
    (38,0.027)(40,0.023)(42,0.019)(44,0.015)(46,0.012)(48,0.009)(50,0.007)
    (52,0.005)(54,0.003)
  };
  \addlegendentry{$df=30$}

  % 说明文字
  \node[font=\scriptsize, align=left] at (axis cs:38,0.18)
    {$df$ 增大 $\Rightarrow$\\趋近正态};
\end{axis}

%% ── 子图标题（外置为节点，与整体标题层次分明）────────────────────
\node[font=\small\bfseries] at ($(ax1.north)+(0,0.3cm)$)
  {$\bar{X}$ 的抽样分布};
\node[font=\small\bfseries] at ($(ax2.north)+(0,0.3cm)$)
  {$(n{-}1)S^2/\sigma^2$ 的 $\chi^2$ 分布};

%% ── 整体标题 ────────────────────────────────────────────────────
\node[font=\large\bfseries]
  at ($(ax1.north)!0.5!(ax2.north)+(0,1.5cm)$)
  {样本均值与样本方差的抽样分布};

\end{tikzpicture}
\end{document}
